% ======================================================================
% 00-opening.tex -- SPEC 0055, frames 1-3.
%
% Insight-first opening for Trackastra / ML lab experts. Three frames:
% title; the questions and the evidence standard; five unexpected findings.
% The proposal map, the "bar", the exploration table, and the "three claims"
% slide are deliberately removed. No tutorial exposition.
%
% Shared lane definitions:
%   - palette: talkred, talkorange, talkteal, talkgray, talklight,
%     talkblue, talkgreen
%   - \talksectiondivider{<large title>}{<one-line subtitle>}
% ======================================================================

% ----------------------------------------------------------------------
% Shared palette (used by all slide lanes).
% ----------------------------------------------------------------------
\definecolor{talkred}{HTML}{C00000}
\definecolor{talkorange}{HTML}{D97706}
\definecolor{talkteal}{HTML}{00756F}
\definecolor{talkgray}{HTML}{5A5A5A}
\definecolor{talklight}{HTML}{EEF2F6}
\definecolor{talkblue}{HTML}{00008C}
\definecolor{talkgreen}{HTML}{1F7A5C}

% ----------------------------------------------------------------------
% \scite{<keys>} -- small inline citation.
% Keeps the reference marker a superscript so citations stay in the text
% flow; the deck does not use bottom captions or footnotes for sources.
% ----------------------------------------------------------------------
\newcommand{\scite}[1]{\textsuperscript{\cite{#1}}}

% ----------------------------------------------------------------------
% \talksectiondivider{<large title>}{<one-line subtitle>}
% One centered divider frame with vertical fill above and below. Section
% lanes call this at top level; #1 uses the TUD slide-title face.
% ----------------------------------------------------------------------
\newcommand{\talksectiondivider}[2]{%
  \begin{frame}
    \vfill
    \centering
    {\TudSlideTitleFont\fontsize{6.5mm}{7.4mm}\selectfont
      \textcolor{tudblue}{#1}\par}
    \vspace{3mm}
    {\normalsize\textcolor{talkgray}{#2}\par}
    \vfill
  \end{frame}
}

% ----------------------------------------------------------------------
% \talktitle{<title text>}
% Wrap slide titles in a 0.90-textwidth parbox so long titles wrap inside
% the frame instead of running past the right edge. \smash gives the title
% zero height, so a wrapped multi-line title does not shift the body into
% the title band; the theme still draws the title overlay and footer.
% ----------------------------------------------------------------------
\newcommand{\talktitle}[1]{\frametitle{\smash{\parbox[t]{0.90\textwidth}{\raggedright\fontsize{4.8mm}{5.3mm}\selectfont #1}}}}

% ----------------------------------------------------------------------
% Frame 1 -- title page.
% ----------------------------------------------------------------------
\begin{frame}[plain]
  \TudTitlePage{Efficient sparse attention and\\self-supervised learned features for\\transformer-based cell tracking}{Leonard Starke \\ Research Assignment \\ 28.10.2026}
\end{frame}

% ----------------------------------------------------------------------
% Frame 2 -- questions and the evidence standard.
% Claim: each question has a pre-committed test, not a post-hoc narrative.
% ----------------------------------------------------------------------
\begin{frame}
  \talktitle{What we set out to learn}
  \centering
  \hyphenpenalty=10000
  \exhyphenpenalty=10000
  \small
  \renewcommand{\arraystretch}{1.45}
  \begin{tabular}{@{}>{\raggedright\arraybackslash}p{61mm}>{\raggedright\arraybackslash}p{64mm}@{}}
    \toprule
    {\raggedright\textbf{Question}} & {\raggedright\textbf{What would count as evidence}}\\
    \midrule
    {\raggedright Where does the cost of the cutoff mask actually live?} &
    {\raggedright Kernel-verified dispatch; production batch $B{=}64$.}\\
    {\raggedright Can a warp manufacture association labels without annotations?} &
    {\raggedright Held-out \texttt{trpL}; pre-registered rule; compute-matched control.}\\
    {\raggedright Which signal carries association: position and geometry, or learned features?} &
    {\raggedright Coordinates-only candidates; frame-CV and LOCO.}\\
    \bottomrule
  \end{tabular}
  \vspace{4mm}

  {\footnotesize The talk reports what changed our mind, not what the proposal hoped.\par}
\end{frame}

% ----------------------------------------------------------------------
% Frame 3 -- five unexpected findings.
% Claim: the talk's payload is these five results, each backed by a measured
% artifact; the rest of the deck supplies the evidence and the bounds.
% ----------------------------------------------------------------------
\begin{frame}
  \talktitle{Five things we did not expect}
  \centering
  \hyphenpenalty=10000
  \exhyphenpenalty=10000
  \small
  \renewcommand{\arraystretch}{1.35}
  \begin{tabular}{@{}p{6mm}>{\raggedright\arraybackslash}p{103mm}p{17mm}@{}}
    \toprule
    \textbf{\#} & {\raggedright\textbf{Finding}} & \textbf{Section}\\
    \midrule
    1 & The mask already reaches a fused cuDNN kernel; the mask itself is the cost. & \S4.1\\
    2 & Removing the cutoff costs accuracy in 12/12 fold--seed cells. & \S4.1\\
    3 & Attention is 10.05\% of the step; the norm chain is 38.7\%. & \S4.1\\
    4 & The warp pretext collapses to all-negative for a structural reason. & \S4.2\\
    5 & Position and geometry beat learned features by 0.015 in LOCO. & \S4.3/4.4\\
    \bottomrule
  \end{tabular}
  \vspace{4mm}

  {\footnotesize Each item is backed by a measured result; open items are named as open.\par}
\end{frame}
